\section{Open Problems} \label{sec-conclusion}

\heading{Randomized Schnorr.} Section~\ref{sec-derand} proves adaptive security, with no query bound in the scheme, for derandomized Schnorr, whose signer derives its nonce from the key and message as in EdDSA. For Schnorr with a randomized signer the situation is unchanged: the trigger branch produces one signature per slot, a randomized signer produces fresh signatures for repeated messages, and the salt of Section~\ref{sec-adaptive} is our only mechanism for reconciling the two. Adaptive security for randomized Schnorr without a salt and without a query bound remains open, as does removing the digest guessing, which is the remaining gap between Corollary~\ref{cor-der-static} and a fully polynomial adaptive result.

\heading{Polynomial security.} Theorem~\ref{thm-adaptive} guesses the digest of one message in each signing hybrid and guesses $(\dig(m^*),\salt^*)$ when treating a trigger-branch forgery. The lossy pre-hash bounds each guess by the image size $r$, so the losses are quasi-polynomial for quasi-polynomial $r$, but they are not polynomial. A proof that punctures at adaptively selected inputs without guessing would remove the remaining loss. Prefix-guessing techniques for adaptively secure signatures~\cite{CCS:RamWat14} and adaptive security analyses for GGM-based puncturable PRFs~\cite{AC:FKPR14,C:JKKKPW17} are natural starting points.

\heading{One hash function for all users.} In our unforgeability constructions the hash key depends on the user's public key $X$, so each user has a different hash function. It is open to construct one hash function that works for all users.

\heading{Beyond lossiness.} Proposition~\ref{prop-lossy-forge} limits what the technique of Section~\ref{sec-lf-kr} can reach: no hash function with a small enumerable image gives unforgeability, so key recovery is the most that lossiness alone can give. Whether some intermediate structure, weaker than the obfuscated conversion function of Section~\ref{sec-io-uf} but stronger than lossiness, suffices for unforgeability is open.

\heading{Weaker assumptions.} Corollary~\ref{cor-bk} uses a strongly efficiently enumerable ELF built from exponential DDH. It would be interesting to construct a post-quantum tunable lossy function with a comparable universal injective mode. For unforgeability, the main question is whether iO can be replaced by a weaker tool.

\heading{Dilithium.} Part~II proves the key-recovery and static results; the salted adaptive transfer is laid out at the end of Section~\ref{sec-dil-static} and remains to be written in full. Independent of the ports, four questions are open: removing the flooding condition from the static result, so that a polynomial modulus suffices; balanced short decompositions at NTT-friendly prime moduli; the sparse challenge set $B_\tau$; and compression with hints. The proof architecture appears compatible with a module setting, but the algebraic reduction and its parameters remain to be carried out: invertibility, zero divisors, the gadget, challenge multiplication, and norm accounting all change.

\heading{Extensions.} It would also be interesting to extend the proofs to key-prefixed Schnorr, multi-user security, and threshold signing. Each setting requires a separate treatment of the signing simulation.
